\newglossaryentry{term-inference}{
    name={Инференс},
    description={процесс получения предсказаний от обученной модели машинного обучения по новым входным данным без изменения ее параметров}
}
\newglossaryentry{term-cold-start}{
    name={Холодный старт},
    description={состояние, при котором инстанс сервиса инференса создается с нуля и перед обработкой первого запроса должен выполнить полную инициализацию, включая загрузку файлов модели}
}
\newglossaryentry{term-weights}{
    name={Веса модели},
    description={набор числовых параметров нейросети, определяющих ее поведение и загружаемых в оперативную и видеопамять перед инференсом}
}
\newglossaryentry{term-chunk}{
    name={Чанк},
    description={фрагмент файла весов фиксированного размера, используемый как минимальная единица кэширования, репликации и P2P-передачи}
}
\newglossaryentry{term-serverless}{
    name={Бессерверные вычисления},
    description={модель развертывания, при которой платформа автоматически управляет выделением и освобождением вычислительных ресурсов в зависимости от нагрузки, включая масштабирование до нуля}
}
\newglossaryentry{term-pod}{
    name={Pod},
    description={минимальная единица развертывания в Kubernetes, объединяющая один или несколько контейнеров с общим сетевым пространством имен и томами хранения}
}
\newglossaryentry{term-daemonset}{
    name={DaemonSet},
    description={ресурс Kubernetes, обеспечивающий запуск экземпляра служебного пода на каждой подходящей worker-ноде кластера}
}
\newglossaryentry{term-scale-to-zero}{
    name={Scale-to-zero},
    description={способность платформы снижать число активных реплик сервиса до нуля при отсутствии входящих запросов}
}
\newglossaryentry{term-control-plane}{
    name={Control Plane},
    description={плоскость управления, в рамках работы представленная Redis-реестром метаданных о чанках и состоянии агентов}
}
\newglossaryentry{term-data-plane}{
    name={Data Plane},
    description={плоскость данных, включающая локальные диски worker-нод, каналы доставки чанков и компоненты, непосредственно обслуживающие чтение весов модели}
}
\newglossaryentry{term-metadata-registry}{
    name={Реестр метаданных},
    description={служба, хранящая сведения о расположении чанков, их размере, контрольных суммах и доступных узлах-держателях}
}
\newglossaryentry{term-fuse}{
    name={Файловая система в пространстве пользователя (FUSE)},
    description={механизм Linux, позволяющий реализовать файловую систему как пользовательский процесс с обработкой обращений через VFS}
}
\newglossaryentry{term-uds}{
    name={Сокет домена Unix},
    description={механизм межпроцессного взаимодействия внутри одного узла, использующий файловую систему вместо IP-стека}
}
\newglossaryentry{term-nfs}{
    name={Сетевая файловая система},
    description={протокол распределенного доступа к удаленным файлам, позволяющий монтировать удаленный каталог как локальный}
}
\newglossaryentry{term-flock}{
    name={\texttt{flock}},
    description={системный вызов для установки рекомендательных файловых блокировок, используемый для координации параллельных загрузок}
}
\newglossaryentry{term-lease}{
    name={Lease},
    description={ресурс Kubernetes Coordination API, представляющий временной токен владения и применяемый для распределенной координации}
}
\newglossaryentry{term-leader-election}{
    name={Leader Election},
    description={алгоритм выбора единственного активного исполнителя привилегированной операции среди конкурирующих процессов}
}
\newglossaryentry{term-cache-hit}{
    name={Cache hit},
    description={ситуация, при которой запрошенные данные найдены в кэше и возвращаются без обращения к более медленному источнику}
}
\newglossaryentry{term-cache-miss}{
    name={Cache miss},
    description={ситуация, при которой данные отсутствуют в текущем уровне кэша и требуется переход к следующему источнику}
}
\newglossaryentry{term-fallback}{
    name={Fallback},
    description={механизм переключения на резервный источник данных при отсутствии или недоступности основного пути чтения}
}
\newglossaryentry{term-spof}{
    name={Единственная точка отказа},
    description={компонент системы, отказ которого приводит к полной недоступности критически важной функциональности}
}
\newglossaryentry{term-peer-discovery}{
    name={Обнаружение узлов},
    description={процесс получения актуального списка агентов, которые могут обслуживать запросы на чтение конкретных чанков модели}
}
\newglossaryentry{term-prefetching}{
    name={Предзагрузка},
    description={заблаговременная загрузка следующих чанков на основе наблюдаемого паттерна чтения для сокращения задержки последующих запросов}
}
